\section{Theoretical Foundation: Bounded Dynamics and Triple Equivalence}
\label{sec:foundation}

\subsection{The Axiom of Boundedness}

We begin with a single foundational axiom:

\begin{axiom}[Boundedness]
\label{ax:bounded}
Physical systems occupy finite domains in phase space.
\end{axiom}

This is not a simplifying assumption but an observational fact. A gas molecule confined to a container cannot escape to infinity. An electron bound to an atom occupies finite spatial extent. An electromagnetic field in a cavity has finite energy. A molecular vibration is constrained by chemical bond lengths. Any system with infinite extent or infinite energy is unphysical.

The implications of boundedness are profound and far-reaching. We prove that boundedness alone necessitates the complete structure of oscillatory dynamics, categorical states, and temporal partitioning.

\subsection{Bounded Dynamics Implies Oscillation}

\begin{proposition}[Oscillatory Necessity]
\label{prop:oscillation}
Any bounded dynamical system with continuous time evolution exhibits oscillatory or quasi-periodic behavior.
\end{proposition}

\begin{proof}
Let the system occupy domain $\mathcal{D} \subset \mathbb{R}^{2n}$ in phase space (for $n$ degrees of freedom), with $\mathcal{D}$ compact. Let $\mathbf{x}(t)$ denote the trajectory starting from $\mathbf{x}_0 \in \mathcal{D}$.

Since $\mathcal{D}$ is bounded, the trajectory cannot escape to infinity. For continuous dynamics, the trajectory is a continuous curve in $\mathcal{D}$. By the Poincaré recurrence theorem \citep{poincare1890probleme}, for measure-preserving flows on bounded domains, almost every trajectory returns arbitrarily close to its initial state infinitely often:
\begin{equation}
\forall \epsilon > 0, \, \exists T(\epsilon) : \, |\mathbf{x}(T) - \mathbf{x}_0| < \epsilon
\end{equation}

For systems with regular (non-chaotic) dynamics, trajectories lie on invariant tori by the KAM theorem \citep{kolmogorov1954conservation,arnold1963small,moser1962invariant}. Motion on a torus is quasi-periodic: a sum of periodic motions with incommensurable frequencies. This is oscillatory behavior.

For chaotic systems, trajectories densely fill accessible phase space but remain confined to $\mathcal{D}$. The density implies repeated encounters with any neighborhood, which is quasi-periodic behavior in the sense of Poincaré recurrence.

In both cases, boundedness implies return, which implies oscillation (periodic or quasi-periodic).
\end{proof}

\subsection{Oscillation Defines Categories}

An oscillating system traverses distinct states. At each instant, it occupies a specific point in phase space. The collection of points visited during evolution defines the oscillation's structure.

\begin{definition}[Category]
\label{def:category}
A \textit{category} is a distinguishable region of phase space that the system's trajectory visits.
\end{definition}

Categories are not imposed externally but emerge from the dynamics itself. Consider a pendulum: at maximum displacement left, maximum displacement right, and zero displacement (center) constitute three distinct categories. The oscillation is the traversal between these categories.

\begin{proposition}[Categorical Structure]
\label{prop:categories}
Oscillatory dynamics and categorical structure are equivalent: oscillation is movement through categories, and categories are the states that oscillation distinguishes.
\end{proposition}

\begin{proof}
Oscillation requires distinguishable states (otherwise there is no periodic structure to identify). Conversely, distinguishable states require dynamics to distinguish them (a static system has no categorical structure). The oscillation period $T$ is the time to traverse all categories and return to the initial category. The number of categories $M$ determines the oscillation's complexity. These are two descriptions of identical mathematical structure.
\end{proof}

\subsection{Categories Partition the Period}

The oscillation period $T$ naturally decomposes into temporal segments corresponding to time spent in each category.

\begin{definition}[Partition]
\label{def:partition}
A \textit{partition} is a temporal segment of the period corresponding to occupation of one categorical state.
\end{definition}

For $M$ categories, the period partitions as:
\begin{equation}
T = \sum_{i=1}^{M} \tau_i
\end{equation}
where $\tau_i$ is the duration of partition $i$.

\begin{proposition}[Partition Structure]
\label{prop:partitions}
The partition structure of an oscillation is equivalent to its categorical structure: each partition corresponds to one category, and the union of all partitions equals the period.
\end{proposition}

\subsection{The Triple Equivalence Theorem}

We now establish the central result of this section:

\begin{theorem}[Triple Equivalence]
\label{thm:triple}
For any bounded dynamical system, the following three descriptions are mathematically equivalent:
\begin{enumerate}
\item \textbf{Oscillatory:} The system exhibits periodic or quasi-periodic motion with characteristic frequencies $\{\omega_i\}$.
\item \textbf{Categorical:} The system traverses $M$ distinguishable states per fundamental period.
\item \textbf{Partition:} The fundamental period $T$ is partitioned into $M$ temporal segments $\{\tau_i\}$.
\end{enumerate}
These are not three separate phenomena but three perspectives on a single underlying structure.
\end{theorem}

\begin{proof}
From Proposition~\ref{prop:oscillation}, boundedness implies oscillation. From Proposition~\ref{prop:categories}, oscillation defines categories. From Proposition~\ref{prop:partitions}, categories partition the period. The three descriptions are logically equivalent.

Quantitatively, the relationships are deterministic:
\begin{align}
\text{Oscillatory frequency:} \quad & \omega = \frac{2\pi}{T} \\
\text{Category traversal rate:} \quad & \frac{dM}{dt} = \frac{M}{T} = \frac{M\omega}{2\pi} \\
\text{Average partition duration:} \quad & \langle \tau \rangle = \frac{T}{M} = \frac{2\pi}{M\omega}
\end{align}

These three quantities are related by:
\begin{equation}
\boxed{\frac{dM}{dt} = \frac{\omega}{2\pi/M} = \frac{1}{\langle \tau \rangle}}
\label{eq:triple_identity}
\end{equation}

Given any one quantity, the others are uniquely determined. The three perspectives are thus informationally equivalent.
\end{proof}

\subsection{Entropy in Three Equivalent Forms}

The triple equivalence implies that thermodynamic entropy admits three equivalent formulations.

\begin{theorem}[Entropy Equivalence]
\label{thm:entropy}
For a bounded system in thermal equilibrium, entropy has three equivalent expressions:
\begin{align}
S_{\text{cat}} &= k_B M \ln n \label{eq:entropy_cat} \\
S_{\text{osc}} &= k_B \sum_i \ln\left(\frac{A_i}{A_0}\right) \label{eq:entropy_osc} \\
S_{\text{part}} &= k_B \sum_a \ln\left(\frac{1}{s_a}\right) \label{eq:entropy_part}
\end{align}
where $M$ is the number of categorical dimensions, $n$ is states per dimension, $A_i$ are oscillation amplitudes, $A_0$ is the ground state amplitude, $s_a$ are partition selectivities, and $k_B$ is Boltzmann's constant.
\end{theorem}

\begin{proof}
\textbf{Categorical entropy} (Equation~\ref{eq:entropy_cat}): A system with $M$ independent categorical dimensions, each admitting $n$ distinguishable states, has total state count $\Omega = n^M$. The Boltzmann entropy formula gives:
\begin{equation}
S = k_B \ln \Omega = k_B \ln(n^M) = k_B M \ln n
\end{equation}

\textbf{Oscillatory entropy} (Equation~\ref{eq:entropy_osc}): Each oscillatory mode with amplitude $A_i$ corresponds to energy $E_i \propto A_i^2$. The number of quantum states accessible at this energy is $\Omega_i \propto A_i/A_0$ (where $A_0$ is the zero-point amplitude). The total entropy is:
\begin{equation}
S = k_B \sum_i \ln \Omega_i = k_B \sum_i \ln\left(\frac{A_i}{A_0}\right)
\end{equation}

\textbf{Partition entropy} (Equation~\ref{eq:entropy_part}): Each partition $a$ with selectivity $s_a$ (probability of transitioning into that partition) corresponds to information content $-\ln s_a$. Summing over all partitions:
\begin{equation}
S = k_B \sum_a \ln\left(\frac{1}{s_a}\right)
\end{equation}

The equivalence follows from the triple equivalence theorem: $M$ categories correspond to $M$ oscillatory modes with equal amplitude spacing, which partition into $M$ segments with equal selectivity. Setting $n = A_i/A_0 = 1/s_a$ for uniform distributions yields:
\begin{equation}
S_{\text{cat}} = k_B M \ln n = k_B \sum_{i=1}^{M} \ln n = S_{\text{osc}} = S_{\text{part}}
\end{equation}
\end{proof}

\subsection{Three-Dimensional Entropy Coordinate Space}

For a general bounded system, complete state specification requires three coordinates corresponding to the three perspectives of the triple equivalence.

\begin{definition}[Entropy Coordinate Space]
\label{def:entropy_space}
The entropy coordinate space is $\Sspace = [0,1]^3$ with coordinates:
\begin{align}
\Sk &= \frac{M_{\text{actual}}}{M_{\text{max}}} \in [0,1] \quad \text{(knowledge: categorical occupation)} \\
\St &= \frac{\phi}{2\pi} \in [0,1] \quad \text{(temporal: oscillatory phase)} \\
\Se &= \frac{t}{T_{\text{rec}}} \in [0,1] \quad \text{(evolution: partition history)}
\end{align}
where $M_{\text{actual}}$ is the number of occupied categories, $M_{\text{max}}$ is maximum possible categories, $\phi$ is oscillatory phase, and $T_{\text{rec}}$ is Poincaré recurrence time.
\end{definition}

This space is bounded by construction—all coordinates lie in $[0,1]$—and three-dimensional because the triple equivalence provides exactly three independent specifications of system state.

\begin{theorem}[Completeness of Entropy Coordinates]
\label{thm:completeness}
The entropy coordinates $(\Sk, \St, \Se)$ provide complete specification of a bounded system's macroscopic state, up to microscopic details indistinguishable at the given resolution.
\end{theorem}

\begin{proof}
Any bounded system's state is characterized by:
\begin{enumerate}
\item Which categories are occupied (knowledge $\Sk$)
\item What phase within the oscillation cycle (temporal $\St$)
\item How the current state was reached (evolution $\Se$)
\end{enumerate}

These three specifications are independent: knowing which categories are occupied does not determine the phase or history; knowing the phase does not determine which categories are accessible or how they were reached; knowing the history does not uniquely determine current occupation or phase.

Conversely, these three specifications are sufficient: given $(\Sk, \St, \Se)$, the system's macroscopic state is determined up to microscopic details that do not affect categorical structure. Different microstates with identical $(\Sk, \St, \Se)$ are indistinguishable at the categorical level.

The space is three-dimensional because three independent coordinates are necessary and sufficient for complete macroscopic state specification.
\end{proof}

\subsection{Fundamental Constants and Scaling}

The triple equivalence introduces natural scales relating categorical counting to physical quantities.

\begin{proposition}[Boltzmann Constant Interpretation]
\label{prop:boltzmann}
Boltzmann's constant $k_B$ converts between categorical rate (categories per second) and energy (joules):
\begin{equation}
k_B = \frac{\hbar \omega}{2\pi T}
\end{equation}
where $\hbar$ is reduced Planck's constant, $\omega$ is oscillation frequency, and $T$ is temperature.
\end{proposition}

\begin{proof}
From the categorical rate equation:
\begin{equation}
\frac{dM}{dt} = \frac{\omega}{2\pi}
\end{equation}

For a quantum harmonic oscillator, energy is $E = \hbar\omega$. Temperature relates to categorical rate by:
\begin{equation}
T = \frac{\hbar}{k_B} \frac{dM}{dt} = \frac{\hbar \omega}{2\pi k_B}
\end{equation}

Solving for $k_B$:
\begin{equation}
k_B = \frac{\hbar \omega}{2\pi T}
\end{equation}

This shows $k_B$ is the conversion factor between energy ($\hbar\omega$) and categorical thermal units ($2\pi T$).
\end{proof}

\subsection{Summary: From Boundedness to Entropy Coordinates}

We have proven that the single axiom of boundedness (Axiom~\ref{ax:bounded}) necessitates:
\begin{enumerate}
\item Oscillatory dynamics (Proposition~\ref{prop:oscillation})
\item Categorical state structure (Proposition~\ref{prop:categories})
\item Temporal partition structure (Proposition~\ref{prop:partitions})
\item Triple equivalence of these three perspectives (Theorem~\ref{thm:triple})
\item Three equivalent entropy formulations (Theorem~\ref{thm:entropy})
\item Three-dimensional entropy coordinate space (Definition~\ref{def:entropy_space})
\item Completeness of entropy coordinates for state specification (Theorem~\ref{thm:completeness})
\end{enumerate}

All subsequent results follow from this foundation. Measurement theory, ternary encoding, temporal resolution, and spectroscopic analysis are consequences of bounded oscillatory dynamics in three-dimensional entropy coordinate space.
